
% --------------------------------------------------------------
\chapter{State of the Art}
\label{ch:state-of-the-art}
% --------------------------------------------------------------

The problem of obtaining the numerical parameters of a graphical
decision model has been approached from several angles. We organise
the literature along four complementary lines and then position the
present work within them.

\section{Statistical learning from data (MLE / EM)}

The classical route to parameterising any graphical model is
maximum-likelihood estimation. When the data is fully observed the CPTs reduce to relative frequencies; with hidden variables or missing
entries the Expectation$-$Maximisation algorithm is the standard
choice. In the specific context of IDs, the work of Ezawa and Gupta~\cite{ezawa1997learning} combined EM for the CPTs with a
greedy search for the causal structure, and extended the procedure to
the utility node via regression on observed outcomes.

The approach is already implemented but requires a sizeable corpus of labelled
historical cases, which in many real applications is simply not available.

\section{Direct expert elicitation}

When historical data is scarce, the established alternative is the
structured elicitation of probabilities and utilities through
techniques such as reference lotteries, indifference judgements and
canonical models (Noisy-OR, Noisy-AND, additive utilities)
\cite{bielza2010modeling}. The literature on multi-criteria decision
making, including the work of Shahzad on the integration of
Bayesian networks and MCDM~\cite{shahzad2022mcdm}, has produced a comprehensive catalogue of methods to reduce the cognitive load on the
expert. 

Even with these aids, however, the resulting judgements are known to be systematically biased: the experimental psychology of
probability assessment has identified three robust distortions —
\emph{availability}, \emph{representativeness} and
\emph{anchoring} — that all degrade the quality of expert-supplied
numbers regardless of the analyst's
training~\cite{bielza2010modeling}. 

The bottleneck therefore
persists on two fronts at once: every CPT row and every utility
entry must be confirmed by a human, the size of the table grows
exponentially with the number of parents, and what is finally
written down is itself an unreliable estimate.

\section{Evolutionary algorithms and metaheuristics}

A third line of work treats parameter elicitation as an
\emph{inverse optimisation} problem: rather than asking the expert
for every parameter, one searches the parameter space for a
configuration that reproduces a weak supervisory signal. Genetic
Algorithms have been used inside the UPM CIG group for this very
purpose~\cite{henriquez2024tfm}, and swarm-intelligence techniques
(PSO, ACO) have been applied to the related problem of learning
unknown behaviours in Interactive Dynamic IDs by Pan et
al.~\cite{pan2025iddids}. Within evolutionary computation, the
family of \emph{Estimation of Distribution Algorithms}
(EDAs)~\cite{larranaga2002eda} has matured into a powerful framework
that, instead of recombining individuals via crossover and mutation,
explicitly learns a probabilistic model of the promising regions of
the search space at every generation.

Recent surveys
\cite{larranaga2024edas} report consistent improvements over
classical GAs in continuous, high-dimensional and constrained
problems, and the open-source library
\textsc{EDAspy}~\cite{soloviev2024edaspy} provides
high-quality implementations of several EDAs, including UMDA, EGNA and KEDA.

Evolutionary algorithms, which we use as estimation methods, allow incorporating knowledge from probabilistic models (parametric and non-parametric) and from the specific decision problem domain into the solution, via constraints.

\section{Generative and LLM-based elicitation}

The most recent trend exploits large language models as artificial experts. Prompted with the structure of the network
and natural-language descriptions of the variables, models such as
GPT-4o, Gemini Pro and Claude 3.5 have been shown to produce CPTs
whose Kullback–Leibler divergence with respect to expert-elicited
tables is comparable to the noise observed across different human
experts~\cite{nafar2025llm4bn}. Closely related, Pan et al.~use
Variational Auto-encoders to generate plausible models of the
unobserved agents in Interactive Dynamic IDs~\cite{pan2024vae}.


These methods are still in their infancy but already suggest powerful
hybrid pipelines in which a generative model produces a promising initial
configuration that is subsequently refined by a metaheuristic.

\section{Beyond exact evaluation}

The Shachter algorithm computes the optimal policy exactly, but only
for regular discrete IDs like the one used in this work. When the
diagram grows too large, mixes continuous and discrete variables, or
carries constraints that the regular formalism cannot express, exact
evaluation is no longer feasible. Two families of methods step in
here: \emph{Decision Programming}~\cite{salo2022decisionprog}, which
rewrites the diagram as an integer optimisation problem and hands it
to a standard solver, and \emph{reinforcement
learning}~\cite{gomez2024rl}, which approximates the optimal policy by
trial and error on a simulator of the diagram. 

Both pursue the same
goal — given the parameters, find the best decisions — which is the
mirror image of the problem studied here, where the decisions (the
expert rules) are given and the parameters are what we look for.
These approaches fall outside the scope of this thesis, but they are
closely related and are worth keeping in mind as the natural
counterpart to the parameter-recovery problem we address.


\section{The inverse problem: rules from CPTs}

A complementary line of work, also originated at UPM, approaches the \emph{inverse} of the problem studied in this thesis. Fern{\'a}ndez del Pozo, Bielza and Lucas~\cite{fdezdelpozo2010cluster} propose the
\textsc{Cluster-KBM2L} methodology, which starts from a fully
specified CPT (more than 8\,000 entries in the case of their
non-Hodgkin lymphoma ID) and extracts \emph{symbolic prognostic
profiles} — qualitative rules of the form ``high-risk patient with
favorable chemotherapy response'' — by clustering the conditional
distributions and synthesizing interpretable descriptions for each
cluster. 

Both works thus address the same interface between the
numerical and the symbolic representation of an ID, but from
opposite directions: \cite{fdezdelpozo2010cluster} goes from CPTs
to rules, whereas this thesis goes from rules to CPTs.

\section{Research gap and positioning of this work}

The review summarized above leads to a clear conclusion:
no published work combines \emph{modern EDAs} with the
\emph{reconstruction of ID parameters from a small set of expert
decision rules}. The closest precedent is the genetic-algorithm
approach of Henr{\'\i}quez~\cite{henriquez2024tfm}. The contribution
of the present thesis is to replace the GA with an EDA that
\emph{explicitly} models the dependency structure of the parameter
space, and to study empirically the impact of this richer
probabilistic model on convergence and final accuracy.
